\section{Τεχνικό Υπόβαθρο}
\label{Background}

Στο κεφάλαιο αυτό, θα χτιστεί το τεχνικό υπόβαθρο που χρειάζεται για την κατανόηση της εργασίας. Συμπεριλαμβάνονται λεπτομέρειες τόσο για έννοιες μηχανικής μάθησης όσο και για βασικές ιδέες της γλωσσολογίας.

Να σημειωθεί ότι η παρουσίαση που ακολουθεί δεν είναι εξαντλητική. Το προτεινόμενο σύστημα αξιοποιεί πολλά δομικά στοιχεία της βαθιάς μάθησης τα οποία αποτελούν πλέον διαδεδομένες και καθιερωμένες πρακτικές, όπως τα συνελικτικά νευρωνικά δίκτυα (CNNs), τα αναδρομικά δίκτυα BiLSTM, οι Transformers και ο μηχανισμός του attention, καθώς και τεχνικές όπως το Dropout, το Layer Normalization και ο optimizer AdamW. Τα στοιχεία αυτά εφαρμόζονται στην εργασία με τον συνήθη τρόπο, χωρίς τροποποιήσεις, και περιγράφονται εκτενώς στη βιβλιογραφία. Για τον λόγο αυτό θεωρούνται γνωστά και η αναλυτική παρουσίασή τους παραλείπεται, ώστε να περιοριστεί η έκταση του κειμένου. Αντιθέτως, παρουσιάζονται αναλυτικά οι έννοιες που είναι πιο εξειδικευμένες και στις οποίες στηρίζονται συγκεκριμένες σχεδιαστικές επιλογές της εργασίας, οι οποίες δεν μπορούν να γίνουν πλήρως κατανοητές χωρίς ακριβή ορισμό. Αυτές είναι τα φωνήματα και τα visemes, πάνω στα οποία βασίζεται τόσο ο ορισμός του προβλήματος όσο και η ομαδοποίηση των φωνημάτων για τη δημιουργία των soft targets, το positional encoding, το οποίο παρέχει την πληροφορία της χρονικής σειράς των frames στους μηχανισμούς attention του μοντέλου (στον Transformer sequence encoder και στο cross-modal attention του audiovisual μοντέλου), οι συναρτήσεις κόστους cross-entropy και CTC, ο συνδυασμός των οποίων αποτελεί την αντικειμενική συνάρτηση της εκπαίδευσης, η απόσταση Levenshtein, στην οποία βασίζονται η μετρική TER και η αλγοριθμική πρόβλεψη λέξεων και τα Mel-Spectrograms, τα οποία αποτελούν την αναπαράσταση του ήχου στην είσοδο του audiovisual μοντέλου.

\subsection{Φωνήματα και Visemes}
\label{Phonemes}
Τα φωνήματα αποτελούν το βασικό αντικείμενο μελέτης της συγκεκριμένης εργασίας. Ένα φώνημα ορίζεται ως ένα σύνολο παρόμοιων ήχων ομιλίας το οποίο γίνεται αντιληπτό από ομιλητές μίας γλώσσας ως ένας βασικός ήχος, πρόκειται δηλαδή για τη μικρότερη πιθανή μονάδα φωνής. Συμβολίζονται συνήθως ως ένας χαρακτήρας εντός δύο forward slash. Για παράδειγμα, το φώνημα /b/ εκπροσωπεί τον ήχο στην αρχή της αγγλικής λέξης bat. Κάθε ομιλούμενη γλώσσα αξιοποιεί τα φωνήματα σε διαφορετικό βαθμό, ενώ ίδια γράμματα ή φθόγγοι μπορεί σε διαφορετικές γλώσσες να προφέρονται τελείως διαφορετικά. Τα προβλήματα αυτά προσπαθεί να λύσει το International Phonetic Alphabet (IPA), το οποίο είναι ένα ``αλφάβητο'' από φωνήματα βασισμένο κυρίως στους λατινικούς χαρακτήρες και ομοιόμορφο για κάθε γλώσσα. Να σημειωθεί ότι το πλήρες IPA είναι υπερβολικά περίπλοκο για την ανάλυση που ακολουθεί, καθώς αρκετά φωνήματα συναντώνται πολύ σπάνια (ή καθόλου) στην αγγλική γλώσσα που μελετάται καθ' όλη τη διάρκεια της εργασίας.

`Οντας η ελάχιστη μονάδα ομιλίας, τα φωνήματα έχουν αποτελέσει το αντικείμενο ποικίλων γλωσσολογικών μελετών, πολλές από τις οποίες αφορούν το ``διάβασμα χειλιών'' (lip reading) ή, αλλιώς, την κατανόηση ομιλίας χρησιμοποιώντας αποκλειστικά οπτικά (και όχι ακουστικά) μέσα. Από τις μελέτες αυτές έχει δημιουργηθεί το πρόβλημα των ``visemes'', τα οποία αποτελούν το σύνολο ήχων ομιλίας με πανομοιότυπη όψη στα χείλια (ή γενικότερα στο πρόσωπο). Να σημειωθεί ότι τα σύνολα αυτά δεν περιορίζονται σε μοναδικά φωνήματα: για παράδειγμα, οι φράσεις ``elephant juice'' και ``i love you'' αποτελούν ένα χιουμοριστικό αλλά έγκυρο παράδειγμα viseme. Στα πλαίσια της συγκεκριμένης μελέτης, στο εξής ως visemes θα αναφέρονται τα φωνήματα κατά την έκφανση των οποίων η όψη του ομιλητή είναι όμοια. Για παράδειγμα, τα φωνήματα /m/, /b/ και /p/ αποτελούν ένα παράδειγμα viseme. Τα visemes προσφέρουν πληροφορία για την τοποθεσία του στόματος στην οποία παράγεται ο ήχος, αλλά όχι για τον τρόπο με τον οποίο παράγεται ο ήχος, το οποίο απαιτεί ακουστικά στοιχεία.

Πέραν της ομαδοποίησης τους μέσω των visemes, τα φωνήματα μπορούν να ομαδοποιηθούν και με βάση 3 ακόμα χαρακτηριστικά: τον τόπο, τον τρόπο εκφοράς τους και το αν είναι έμφωνα ή όχι. Για να παρουσιαστούν με μεγαλύτερη ακρίβεια τα παραπάνω, συμπεριλαμβάνεται το επίσημο chart φωνημάτων του IPA. Αγνοώντας τις υπερβολικές λεπτομέρειες στο Σχήμα \ref{fig:IPAChart} και εστιάζοντας την προσοχή στον πρώτο πίνακα με τίτλο CONSONANTS (PULMONIC) και στον αμέσως κάτω δεξιά χάρτη με τίτλο VOWELS, μπορούν να εξαχθούν τα παρακάτω συμπεράσματα.

Παρατηρούμε ότι τα σύμφωνα χωρίζονται σε 11 κατηγορίες όσον αφορά τον τόπο εκφοράς τους: τα δειχιλικά (bilabial, παράγονται χρησιμοποιώντας τα 2 χείλη), τα χειλοδοντικά (labiodental, παράγονται με τη συμμετοχή χειλιών και δοντιών), τα οδοντικά (dental, με τα δόντια), τα φατνιακά (alveolar, παράγονται με τη γλώσσα να ακουμπά ή να πλησιάζει τη φατνιακή ακρολοφία), τα μεταφατνιακά (postalveolar, παράγονται με τη γλώσσα να τοποθετείται πίσω από τη φατνιακή ακρολοφία), τα ανακεκαμμένα (retroflex, παράγονται από την άκρη της γλώσσας όταν κάμπτεται προς τα πίσω προς τον ουρανίσκο), τα ουρανικά (palatals, παράγονται υψώνοντας το κυρίως τμήμα της γλώσσας στον σκληρό ουρανίσκο), τα υπερωικά (velar, όταν το πίσω μέρος της γλώσσας αγγίζει ή πλησιάζει τον μαλακό ουρανίσκο), τα σταφυλικά (uvular, με το πίσω μέρος της γλώσσας να αγγίζει ή να πλησιάζει τη σταφυλή), τα φαρυγγικά (pharyngeal, παράγονται περιορίζοντας τον φάρυγγα) και τέλος τα γλωττιδικά (glottal, χρησιμοποιώντας τη γλωττίδα). Το χαρακτηριστικό αυτό είναι που επηρεάζει κατά κύριο λόγο τα visemes, καθώς ο τρόπος με τον οποίον παράγεται ένα φώνημα και το αν είναι έμφωνο ή όχι πολλές φορές γίνεται κρυφό λόγω της μορφολογίας των χειλιών, που επηρεάζεται προφανώς από τον τόπο εκφοράς του φωνήματος.

Πέραν του τόπου εκφοράς, τα σύμφωνα διαχωρίζονται στον πίνακα με βάση τον τρόπο εκφοράς τους ως εξής: τα έκκροτα/κλειστά σύμφωνα (plosive, παράγονται όταν διακόπτεται τελείως η ροή του αέρα και απελευθερώνεται ξαφνικά), τα ρινικά (nasal, επιτρέπουν στον αέρα να διαφύγει μέσω της ρινικής κοιλότητας και όχι του στόματος), τα μέσω τρίλιας (trill, παράγονται μέσω γρήγορης δόνησης λόγω της ροής αέρα), τα παλμού/κτύπου (tap/flap, όταν η γλώσσα ακουμπάει στιγμιαία ένα σημείο του στόματος), τα τριβόμενα (fricative, όταν ο αέρας περνάει από στενή δίοδο δημιουργώντας τριβή), τα πλευρικά τριβόμενα (lateral fricative, όπου ο αέρας παράγει τριβή καθώς περνάει από τα πλάγια της γλώσσας), τα προσεγγιστικά (approximant, όπου δύο αρθρωτές πλησιάζουν χωρίς να παράγουν τυρβώδη ροή αέρα ή ολικό κλείσιμο), και τα πλευρικά προσεγγιστικά (lateral approximant, μπλοκάροντας το κέντρο της στοματικής κοιλότητας στην περιοχή του αρθρωτή και επιτρέποντας τη ροή του αέρα από τα πλάγια). Πειραματιζόμενοι με τα παραπάνω μπορούμε να συνειδητοποιήσουμε πώς παράγονται τα φωνήματα που χρησιμοποιούμε, όπως και να αρθρώσουμε φωνήματα που ίσως δεν έχουμε συνηθίσει (π.χ. στα ελληνικά δεν υπάρχει πλευρικό τριβόμενο φώνημα, αλλά σε γλώσσες όπως τα ουαλικά υπάρχει).

Τέλος, ο χαρακτηρισμός έμφωνο ή άφωνο φαίνεται από τη θέση των συμφώνων στο εκάστοτε κελί. Αυτά που είναι τοποθετημένα προς τα δεξιά είναι έμφωνα, ενώ αυτά προς τα αριστερά είναι άφωνα. Τα φωνήεντα, από την άλλη είναι όλα έμφωνα και διαχωρίζονται με βάση το ύψος της γλώσσας (στον κατακόρυφο άξονα), την προώθησή της (στον οριζόντιο άξονα) και το αν τα χείλη είναι στρογγυλεμένα (rounded) ή όχι. Όσον αφορά το ύψος, οι κατηγορίες είναι close, close-mid, open-mid και open, ενώ στον οριζόντιο άξονα οι πιθανές επιλογές είναι front, central και back. Όπου στον χάρτη εμφανίζονται δύο φωνήματα, το δεξιότερο παράγεται με στρογγυλεμένα χείλη.

Με γνώση των παραπάνων και αντίστοιχη προσαρμογή στις ανάγκες της εκάστοτε γλώσσας ή, στην περίπτωση περιορισμένου περιβάλλοντος, στις ιδιαιτερότητες του λεξιλογίου, τα φωνήματα μπορούν να ομαδοποιηθούν με τους παραπάνω κανόνες. Έτσι, σε πλαίσια lip reading, είναι δυνατόν να χτιστούν μέθοδοι εφαρμογής διαφορετικής ποινής με βάση την ομοιότητα των φωνημάτων (είτε σε επίπεδο visemes είτε σε επίπεδο χαρακτηριστικών, όπως παρουσιάστηκε) ή και μέθοδοι πρόβλεψης λέξεων που να συνυπολογίζουν την ομοιότητα φωνημάτων κατά την αντικατάστασή τους (για περισσότερα βλ. παράγραφο \ref{PhonemeExtraction}).

\begin{figure}[h]
	\centering
	\includegraphics[width=\textwidth]{IPAChart.png}
	\caption{Επίσημο chart φωνημάτων του IPA, έκδοση 2020}
	\label{fig:IPAChart}
\end{figure}

\subsection{Έννοιες Μηχανικής Μάθησης}
\label{MachineLearningBackground}

Στην ενότητα αυτή ορίζονται οι διάφορες έννοιες της μηχανικής μάθησης και της βαθιάς μάθησης που χρησιμοποιούνται παρακάτω. Ταυτόχρονα, αναπτύσσεται ο συμβολισμός που τηρείται στο υπόλοιπο μέρος της εργασίας.

\paragraph{Positional Encoding}

Στη βαθιά μάθηση, η σειρά των tokens έχει σημαντική επίδραση στην κατανόηση του συνολικού περιεχομένου. Το κλασικό παράδειγμα ``Allen walks dog'' και ``dog walks Allen'' αποδεικνύει τον παραπάνω ισχυρισμό: οι δύο φράσεις έχουν πολύ διαφορετική σημασία, παρά το γεγονός ότι παράγονται από τις ίδιες λέξεις (τα tokens σε αυτήν την περίπτωση). Μοντέλα όπως τα επαναλαμβανόμενα νευρωνικά δίκτυα (RNNs) ή τα μοντέλα μακροπρόθεσμης και βραχυπρόθεσμης μνήμης (long-short term memory, LSTM) διαθέτουν μηχανισμούς που αντιλαμβάνονται και αξιοποιούν τη σειρά των tokens. Αντιθέτως, αρχιτεκτονικές που βασίζονται στους transformers επεξεργάζονται τα tokens εν παραλλήλω και άρα δε διατηρούν πληροφορία για τη σειρά των tokens.

Ο τρόπος με τον οποίον επιλύεται το πρόβλημα αυτό είναι χρησιμοποιώντας Positional Encoding που προσθέτει πληροφορία για την σειρά των tokens στο embedding τους, συγκεκριμένα μέσω Sinusoidal Encoding. Στην εργασία αυτή χρησιμοποείται το sinusoidal encoding όπως παρουσιάζεται στο \cite{vaswani2023attentionneed}, που υπολογίζεται ως εξής:

\[
\mathrm{PE}(k,2i) = \sin\left(k/n^{2i/d}\right)
\]
\[
\mathrm{PE}(k,2i+1) = \cos\left(k/n^{2i/d}\right)
\]

όπου $\mathrm{PE}$ είναι η συνάρτηση positional encoding, $k$ είναι η θέση του token, $i$ είναι η συγκεκριμένη διάσταση, $n$ ένας μεγάλος επιλεγμένος αριθμός (τυπικά 10000) και $d$ η διάσταση του μοντέλου. Προσθέτοντας (ή σε κάποιες περιπτώσεις συνενώνοντας) τους όρους που παράγονται με αυτόν τον κανόνα με τα embeddings των tokens, δίνεται στο μοντέλο η έξτρα πληροφορία της τοποθεσίας του κάθε token.

\paragraph{Συνάρτηση Κόστους Cross-Entropy}

Η συνάρτηση κόστους cross-entropy αποτελεί το τυπικό objective εκπαίδευσης, όταν είναι επιθυμητό ένας ταξινομητής να προσεγγίζει μία στοχευμένη κατανομή. Στη θεωρία πληροφορίας, η cross-entropy μεταξύ δύο κατανομών πιθανότητας $p$ και $q$, στο ίδιο set γεγονότων, αποτυπώνει τον αριθμό των bits που χρειάζονται για να κατηγοριοποιήσουν ένα γεγονός επιλεγμένο από το set όταν η κωδικοποίηση είναι βέλτιστη για την εκτιμώμενη κατανομή πιθανότητας $q$ και όχι την πραγματική $p$. Υπολογίζεται από τον τύπο:

\[
H(p,q) = - \sum_{i}p_i \cdot \log(q_i)
\]

όπου $p_i$ και $q_i$ είναι οι πιθανότητες που αποδίδονται από τις δύο κατανομές για το $i$-οστό γεγονός. Σε περίπτωση που η πραγματική κατανομή $p$ έχει τιμή 1 για μία κλάση και τιμή 0 σε οποιαδήποτε άλλη περίπτωση (αντίστοιχα, σε περίπτωση που τα targets σε ένα πρόβλημα ταξινόμησης ακολουθούν τη σύμβαση one-hot encoding), τότε ο τύπος γίνεται $H(p,q) = -\log(q_i)$, δηλαδή η γνωστή αρνητική λογαριθμική πιθανοφάνεια (negative log-likelihood). Στο πρόβλημα του lip reading, για λόγους που γίνονται κατανοητοί παρακάτω, κρίνεται χρίσημο οι στόχοι να μην είναι απόλυτοι (hard), δηλαδή με πιθανότητα 1 στη σωστή κλάση και 0 αλλού, αλλά εξομαλυμένοι (soft).

Η πρακτική αυτή, κατά την οποία ένα ποσοστό της μάζας πιθανότητας αναδιανέμεται σε άλλες κλάσεις πέραν της σωστής, ονομάζεται label smoothing (εξομάλυνση ετικετών) και έχει ως στόχο να αποτρέψει την υπερβολική αυτοπεποίθηση του μοντέλου και άρα να βελτιώσει τη γενίκευση αυτού. Στη μελέτη αυτή, όπως φαίνεται στην παράγραφο \ref{PhonemeExtraction}, η εξομάλυνση αυτή δεν είναι ομοιόμορφη για όλες τις λανθασμένες κλάσεις αλλά βασίζεται στην ομοιότητα που παρουσιάζουν τα φωνήματα, τόσο οπτικά μέσω των visemes όσο και από το τελικό ηχητικό τους αποτέλεσμα (μέσω των τρόπων, τόπων παραγωγής τους και το αν είναι έμφωνα ή όχι).

Η εφαρμογή της cross-entropy loss στη μελέτη αυτή επιφέρει την ανάγκη δημιουργίας ορθών targets για κάθε frame ενός δεδομένου βίντεο. Επομένως, τα δεδομένα πρέπει να προεπεξεργαστούν καταλλήλως ώστε να αντιστοιχισθεί κάθε frame με το φώνημα που εκφέρεται (ή τη σιγή). Πρόκειται για την παλαιότερη και πιο άμεση πρακτική για την εξαγωγή ετικετών της αλληλουχίας, η οποία χρησιμοποείται από το \cite{GRAVES2005602} και επιτρέπει την επιβλεπόμενη μάθηση σε επίπεδο frame.

\paragraph{Συνάρτηση κόστους CTC}

Για να μην απαιτείται η παραπάνω δημιουργία targets σε επίπεδο frames, εισήχθη από το \cite{inproceedings} η συνάρτηση κόστους CTC (Connectionist Temporal Classification). Η συνάρτηση αυτή δεν απαιτεί label για κάθε frame, αλλά απλά μία αλληλουχία φωνημάτων χωρίς segmentation, εκπαιδεύοντας το μοντέλο να αναθέτει υψηλή πιθανότητα στην αλληλουχία αυτή υπό ενός κανόνα ``κατάρρευσης'' (collapsing rule). Προστίθεται το σύμβολο του κενού και κάθε αλληλουχία σε επίπεδο frame που

\begin{itemize}
\item μετά την αφαίρεση των κενών συμβόλων και
\item μετά την κατάρρευση διαδοχικών ίδιων συμβόλων σε ένα
\end{itemize}
οδηγεί στην αρχική αλληλουχία, θεωρείται σωστή. Έτσι, πολλές διαφορετικές αλληλουχίες φωνημάτων τελικά θεωρούνται σωστές, με αποτέλεσμα η ακριβής πρόβλεψη του χρονισμού του κάθε φωνήματος να απαιτεί πολλά περισσότερα δεδομένα (όπως φαίνεται από την επίδοση του V4 στο \ref{Experiment1}).

Πιο συγκεκριμένα, έστω μία πιθανή ακολουθία με μήκος $T$ που προβλέπει το μοντέλο $\pi = (\pi_1, \pi_2, ..., \pi_T)$, που ονομάζεται μονοπάτι (path). Η συνάρτηση κατάρρευσης $\mathcal{B}(\pi)$ εκτελεί τη λειτουργία που αναφέρθηκε παραπάνω και υπολογίζει από το αρχικό μονοπάτι το collapsed sequence. Υποθέτοντας ανεξαρτησία μεταξύ των χρονικών βημάτων, η πιθανότητα ενός path δεδομένης της ακολουθίας εισόδου $X$ υπολογίζεται ως $P(\pi|X) = \prod_{t=1}^{T} y_{t}^{\pi_t}$, όπου $y_{t}^{\pi_t}$ είναι η πιθανότητα που έδωσε το δίκτυο για το σύμβολο $\pi_t$ στο frame $t$. 

Όπως αναφέρθηκε προηγουμένως, η ίδια τελική ακολουθία μπορεί να παραχθεί από πολλά διαφορετικά αρχικά alignments. Άρα, η πιθανότητα της τελικής ακολουθίας $Y$ ορίζεται ως

\[
P(Y|X) = \sum_{\pi \in \mathcal{B}^{-1}(Y)} P(\pi|X)
\]

και η αντικειμενική συνάρτηση υπολογίζεται ως 

\[
\mathcal{L}_c = -\log P(Y|X) = -\log \left(\sum_{\pi \in \mathcal{B}^{-1}(Y)} P(\pi|X)\right)
\]

Στην πράξη, καθώς το πλήθος των πιθανών paths αυξάνεται εκθετικά με το μήκος της ακολουθίας, δεν είναι εφικτό να απαριθμιστούν όλα από τα παραπάνω μονοπάτια. Αντ' αυτού, χρησιμοποιείται ο αλγόριθμος Forward-Backward, μία μορφή δυναμικού προγραμματισμού, ο οποίος υπολογίζει την πιθανότητα της επιθυμητής αλληλουχίας σε χρόνο $O(T \cdot U)$, όπου $T$ είναι το μήκος της εισόδου και $U$ είναι το μήκος της targeted ακολουθίας. Ο αλγόριθμος αυτός ορίζει forward και backward μεταβλητές που αντιστοιχούν στις πιθανότητες να έχει παραχθεί ένα συγκεκριμένο πρόθεμα ή επίθεμα της ακολουθίας στόχου μέχρι μία χρονική στιγμή. Συνδυάζοντας τις πιθανότητες αυτές υπολογίζεται το συνολικό άθροισμα όλων των έγκυρων αλληλουχιών, χωρίς να εξετάζεται κάθε πιθανό path ξεχωριστά.

Τα πλεονεκτήματα της εκπαίδευσης με την παραπάνω αντικειμενική συνάρτηση είναι η εξάλειψη της απαίτησης ετικετών σε επίπεδο frame, η ικανότητα χρήσης διαφορετικών μηκών εισόδου και εξόδου, η αυτόματη μάθηση της χρονικής ευθυγράμμισης (temporal alignment) των φωνημάτων και ότι επιτρέπει την εκπαίδευση end-to-end. Αποτελεί μειονέκτημα, ωστόσο, η υπόθεση ότι οι προβλέψεις σε διαδοχικά χρονικά βήματα είναι ανεξάρτητες δεδομένης της εισόδου, γεγονός που περιορίζει τη μοντελοποίηση ισχυρών εξαρτήσεων μεταξύ των συμβόλων εξόδου. Άλλα ελαττώματα της χρήσης της συνάρτησης αυτής είναι ότι δεν ενσωματώνεται γλωσσικό μοντέλο κατά την εκπαίδευση, με τη γλωσσική πληροφορία να προστίθεται μόνο κατά την αποκωδικοποίηση, όπως επίσης ότι παρουσιάζονται δυσκολίες κατά την επεξεργασία μεγάλων ακολουθιών ή ακολουθιών με πολλά επαναλαμβανόμενα σύμβολα.

\subsection{Άλλες Χρήσιμες Έννοιες}
\label{OtherBackground}
\paragraph{Απόσταση Levenshtein}

Στη θεωρία πληροφορίας, η απόσταση Levenshtein (ή αλλιώς edit distance) είναι μία μετρική που αφορά την απόσταση δύο αλληλουχιών χαρακτήρων (strings). Αφορά τον ελάχιστο αριθμό αλλαγών μοναδικού χαρακτήρα (αφαιρέσεων, προσθέσεων ή μετατροπών) ώστε ένα string να μετατραπεί σε κάποιο άλλο. Συγκεκριμένα, έστω $a$ και $b$ δύο συμβολοσειρές με μήκη $|a|$ και $|b|$ αντίστοιχα. Η απόσταση Levenshtein των δύο αυτών συμβολοσειρών υπολογίζεται με τον παρακάτω αναδρομικό τύπο ως:

\[
\operatorname{lev}(a,b) = \begin{cases}
|a| & \text{αν } |b| = 0 \\
|b| & \text{αν } |a| = 0 \\
\operatorname{lev}(\operatorname{tail}(a), \operatorname{tail}(b)) & \text{αν } \operatorname{head}(a) = \operatorname{head}(b) \\
1 + \min \begin{cases}  \operatorname{lev}(\operatorname{tail}(a),b) \\
								\operatorname{lev}(a,\operatorname{tail}(b)) \\
								\operatorname{lev}(\operatorname{tail}(a),\operatorname{tail}(b)) \\
			\end{cases} & \text{αλλιώς} \\
			\end{cases}
\]

όπου $\operatorname{tail}(x)$ μίας συμβολοσειράς $x$ είναι η συμβολοσειρά όλων των χαρακτήρων της $x$ δίχως τον πρώτο και $\operatorname{head}(x)$, αντίστοιχα, είναι μόνο ο πρώτος χαρακτήρας της $x$.

Στην πράξη, για να αποφευχθεί η χρήση του παραπάνω αναδρομικού ορισμού, η οποία αυξάνει το υπολογιστικό κόστος εκθετικά με το μήκος των συμβολοσειρών, χρησιμοποιούνται τεχνικές δυναμικού προγραμματισμού για τον υπολογισμό της απόστασης Levenshtein μεταξύ δύο συμβολοσειρών. Παρακάτω δίνεται ένα παράδειγμα υπολογισμού της απόστασης Levenshtein των συμβολοσειρών ``kitten'' και ``sitting''.

Παρατηρούμε ότι ο ελάχιστος αριθμός αλλαγών σε χαρακτήρες που πρέπει να γίνουν για να ταυτιστούν οι δύο συμβολοσειρές είναι 3, και συγκεκριμένα:

\begin{itemize}
\item kitten $\rightarrow$ sitten (μετατροπή του ``k'' σε ``s'')
\item sitten $\rightarrow$ sittin (μετατροπή του ``e'' σε ``i'')
\item sittin $\rightarrow$ sitting (προσθήκη του ``g'')
\end{itemize}

Η μετατροπή δεν μπορεί να γίνει με λιγότερες κινήσεις και άρα η Levenshtein απόσταση των συμβολοσειρών ``kitten'' και ``sitting'' είναι 3.

\paragraph{Mel-Spectrogram}

Τα Mel-Spectrograms αποτελούν μία από τις πλέον διαδεδομένες αναπαραστάσεις ήχου στην επεξεργασία ομιλίας και στη μηχανική μάθηση, καθώς προσφέρουν μια συμπαγή, χρονικά-συχνοτικά δομημένη περιγραφή ενός ηχητικού σήματος που προσεγγίζει τον τρόπο με τον οποίο το ανθρώπινο αυτί αντιλαμβάνεται τον ήχο. Παράγονται υπολογίζοντας το φάσμα του σήματος μέσω του Μετασχηματισμού Fourier Βραχέος Χρόνου (Short-Time Fourier Transform - STFT), μέσω του οποίου το σήμα χωρίζεται σε επικαλυπτόμενα χρονικά παράθυρα, για το καθένα από τα οποία υπολογίζεται το φάσμα ισχύος, παράγοντας μία δισδιάστατη αναπαράσταση συχνοτήτων και χρόνου. Το κλασικό φασματογράφημα, ωστόσο, απεικονίζει τις συχνότητες σε γραμμική κλίμακα, κάτι που δεν αντιστοιχεί με την ψυχοακουστική αντίληψη του ήχου από τον άνθρωπο, καθώς ο άνθρωπος διακρίνει τις χαμηλότερες συχνότητες με πολύ μεγαλύτερη ευκρίνεια από τις υψηλότερες. Έτσι, μέσω μίας αλληλουχίας τριγωνικών φιλτρών, το φάσμα προβάλλεται σε μία λογαριθμική κλίμακα που συμπιέζει τις υψηλές συχνότητες περισσότερο από τις χαμηλές, προσομοιώνοντας την ανθρώπινη ακοή. Μέσω της διαδικασίας αυτής παράγεται το Mel-Spectrogram ενός ηχητικού σήματος, μειώνοντας σημαντικά τη διάσταση (συνήθως χρησιμοποιούνται περίπου 100 μπάντες Mel, ενώ οι αρχικές συνιστώσες του STFT μπορεί να είναι μέχρι και μερικές χιλιάδες) αλλά διατηρώντας σημαντικό μέρος της πληροφορίας που μπορεί να χρησιμοποιηθεί για την ανάκτηση του ήχου (ή, τουλάχιστον, του τμήματος που ακούγεται από τον άνθρωπο).

Πιο συγκεκριμένα, έστω ένα ηχητικό σήμα $y[n]$, δειγματοληπτημένο με ρυθμό $f_s$. Ο υπολογισμός του Mel-Spectrogram γίνεται στα ακόλουθα βήματα:

\begin{enumerate}
\item \textbf{Μετασχηματισμός Fourier Βραχέος Χρόνου.} Το σήμα χωρίζεται σε επικαλυπτόμενα παράθυρα μήκους $N_f$ δειγμάτων, τα κέντρα των οποίων απέχουν μεταξύ τους $N_h$ δείγματα (βήμα ολίσθησης ή hop), και σε κάθε παράθυρο εφαρμόζεται ο Διακριτός Μετασχηματισμός Fourier:
\[
X(t,k) = \sum_{n=0}^{N_f-1} w[n] \, y\!\left[t N_h + n - \tfrac{N_f}{2}\right] e^{-j 2\pi k n / N_f}, \quad k = 0, 1, \ldots, \tfrac{N_f}{2}
\]
όπου $t$ είναι ο δείκτης του χρονικού παραθύρου (frame), $k$ ο δείκτης της συχνοτικής συνιστώσας, η οποία αντιστοιχεί στη συχνότητα $f_k = k f_s / N_f$ σε Hz, και $w[n]$ μία συνάρτηση παραθύρου. Στην εργασία χρησιμοποιείται το παράθυρο Hann, $w[n] = \frac{1}{2}\left(1 - \cos\frac{2\pi n}{N_f}\right)$, ενώ θεωρείται $y[n] = 0$ εκτός του διαστήματος ορισμού του σήματος (zero padding), ώστε το $t$-οστό παράθυρο να είναι κεντραρισμένο στο δείγμα $t N_h$. Λόγω της συμμετρίας του μετασχηματισμού για πραγματικά σήματα, αρκούν οι συνιστώσες $k \leq N_f/2$.

\item \textbf{Φάσμα ισχύος.} Από τον μετασχηματισμό κρατείται μόνο το πλάτος, υψωμένο στο τετράγωνο:
\[
P(t,k) = \lvert X(t,k) \rvert^2
\]

\item \textbf{Κλίμακα mel.} Η κλίμακα mel αντιστοιχίζει τη συχνότητα $f$ (σε Hz) σε μία ψυχοακουστική κλίμακα τονικού ύψους, στην οποία ίσες αποστάσεις γίνονται αντιληπτές από τον άνθρωπο ως ίσες μεταβολές του τόνου \cite{Stevens1937}. Μία συνήθης διατύπωσή της είναι $\operatorname{mel}(f) = 2595 \log_{10}(1 + f/700)$. Στην εργασία χρησιμοποιείται η εκδοχή του Slaney \cite{Slaney1998}, η οποία είναι γραμμική κάτω από το 1 kHz και λογαριθμική πάνω από αυτό:
\[
\operatorname{mel}(f) = \begin{cases}
\dfrac{3f}{200} & f < 1000 \\[2mm]
15 + \dfrac{27 \ln(f/1000)}{\ln 6.4} & f \geq 1000
\end{cases}
\]

\item \textbf{Τριγωνικά φίλτρα.} Ορίζονται $N_m + 2$ συχνότητες $f_0 < f_1 < \ldots < f_{N_m+1}$, ισαπέχουσες στην κλίμακα mel από το $\operatorname{mel}(0)$ έως το $\operatorname{mel}(f_s/2)$. Το $i$-οστό φίλτρο ($i = 1, \ldots, N_m$) έχει κέντρο τη συχνότητα $f_i$ και τριγωνική μορφή:
\[
H_i(f) = \begin{cases}
\dfrac{f - f_{i-1}}{f_i - f_{i-1}} & f_{i-1} \leq f \leq f_i \\[2mm]
\dfrac{f_{i+1} - f}{f_{i+1} - f_i} & f_i \leq f \leq f_{i+1} \\[2mm]
0 & \text{αλλιώς}
\end{cases}
\]
Στην κανονικοποίηση του Slaney, που χρησιμοποιείται και στην εργασία, κάθε φίλτρο πολλαπλασιάζεται επιπλέον με $2/(f_{i+1} - f_{i-1})$, ώστε όλα τα φίλτρα να έχουν το ίδιο εμβαδό. Λόγω της κλίμακας mel, τα φίλτρα είναι στενά στις χαμηλές συχνότητες και όλο και πλατύτερα στις υψηλές.

\item \textbf{Προβολή στα φίλτρα.} Κάθε στοιχείο του Mel-Spectrogram είναι η ενέργεια του φάσματος που ``συλλέγει'' το αντίστοιχο φίλτρο:
\[
S(i,t) = \sum_{k=0}^{N_f/2} H_i(f_k) \, P(t,k), \quad i = 1, \ldots, N_m
\]
Ισοδύναμα, το Mel-Spectrogram προκύπτει ως γινόμενο του πίνακα των φίλτρων (filterbank), διαστάσεων $N_m \times (N_f/2 + 1)$, με τον πίνακα του φάσματος ισχύος.

\item \textbf{Λογαριθμική συμπίεση.} Τέλος, οι τιμές μετατρέπονται σε dB:
\[
S_d(i,t) = 10 \log_{10} \frac{\max\left(S(i,t), \epsilon_0\right)}{\max_{i,t} S(i,t)}
\]
όπου $\epsilon_0$ μία μικρή σταθερά (στην εργασία $10^{-10}$) που αποτρέπει τον λογάριθμο του μηδενός. Καθώς ως αναφορά χρησιμοποιείται η μέγιστη τιμή του Mel-Spectrogram, η μέγιστη τιμή του $S_d$ είναι 0 dB, ενώ στην εργασία τιμές χαμηλότερες από $-80$ dB αποκόπτονται στην τιμή $-80$ dB.
\end{enumerate}

Επομένως, από ένα σήμα μήκους $L$ δειγμάτων προκύπτει ένας πίνακας $S_d \in \mathbb{R}^{N_m \times T}$, με $T = 1 + \lfloor L / N_h \rfloor$ χρονικά frames. Οι παράμετροι $N_f$ και $N_h$ καθορίζουν τον συμβιβασμό ανάμεσα στη συχνοτική και τη χρονική ανάλυση: μεγαλύτερο παράθυρο $N_f$ δίνει καλύτερη συχνοτική ανάλυση (απόσταση $f_s/N_f$ Hz μεταξύ διαδοχικών συνιστωσών) αλλά λιγότερο εντοπισμένη χρονικά πληροφορία, ενώ το βήμα $N_h$ ορίζει τη χρονική απόσταση $N_h/f_s$ μεταξύ διαδοχικών frames. Στην εργασία ο υπολογισμός γίνεται με τη βιβλιοθήκη librosa \cite{McFee2015librosa}, ενώ οι τιμές των παραμέτρων που χρησιμοποιούνται παρουσιάζονται στην ενότητα \ref{MelSpectrogram}.

Η μετατροπή σε κλίμακα dB του τελευταίου βήματος εφαρμόζεται, καθώς και η ένταση του ήχου γίνεται αντιληπτή στον άνθρωπο σε λογαριθμική κλίμακα. Η επιλογή αυτή συμβάλλει, επίσης, στη σταθεροποίηση κατά την εκπαίδευση των νευρωνικών δικτύων. Σε γενικές γραμμές, τα Mel-Spectrograms έχουν υιοθετηθεί ευρέως σε εφαρμογές όπως η αναγνώριση ομιλίας (speech recognition) για τους παρακάτω 3 βασικούς λόγους:

\begin{itemize}
\item μειώνουν τη διάσταση και αυξάνουν την υπολογιστική αποδοτικότητα σε σχέση με την κυματομορφή του ήχου ή το φασματογράφημα αυτής

\item η λογαριθμική κλίμακα ενισχύει τις περιοχές των συχνοτήτων που φέρουν την περισσότερη πληροφορία (όπως τις formants συχνότητες της ανθρώπινης ομιλίας)

\item η δισδιάστατή τους δομή επιτρέπει την εφαρμογή μεθόδων μηχανικής μάθησης εμπνευσμένων από την επεξεργασία εικόνας, όπως των συνελικτικών νευρωνικών δικτύων (CNNs)
\end{itemize}

Παράλληλα, καθώς η μετατροπή από κυματομορφή σε mel-spectrogram είναι ντετερμινιστική συνάρτηση σταθερών υπερπαραμέτρων (ρυθμός δειγματοληψίας, μήκος παραθύρου, βήμα ολίσθησης, αριθμός μπάντων mel), αποτελεί μια αναπαραγώγιμη και εύκολα προ-υπολογίσιμη (cacheable) αναπαράσταση, γεγονός που την καθιστά ιδιαίτερα βολική για χρήση ως είσοδο σε αγωγούς εκπαίδευσης μεγάλης κλίμακας.